\documentclass[compress,10pt,xcolor=dvipsnames]{beamer}
\usepackage{pgfpages}
\setbeameroption{show notes on second screen}
\setbeamerfont{note page}{size=\footnotesize}

\usepackage[english]{babel}
\usepackage[utf8]{inputenc}
\usepackage{times}
\usepackage[T1]{fontenc}
\usepackage{alltt}

\usepackage{tikz}
\usetikzlibrary{calc,arrows,positioning,decorations.markings,shapes}
\usepackage{amsfonts}
\usepackage{amssymb}
\usepackage{amsmath}
\usepackage{graphicx}
\usepackage{tabularx}
\usepackage{hyperref}
\usepackage{natbib}
\setbeamertemplate{bibliography item}[text]

%% COLORS
\definecolor{awesome}{RGB}{3, 149, 222}
\definecolor{high1}{RGB}{3, 149, 222}
\definecolor{high2}{HTML}{F58137}
\definecolor{high3}{RGB}{120, 81, 169}
\definecolor{lightergrey}{rgb}{0.9, 0.9, 0.9}

%% TIKZ NODES & STYLES
\tikzset{
  boxnode/.style={draw=awesome!80, rectangle, rounded corners=2pt, minimum width=2.2cm, minimum height=1.0cm, align=center, thick},
  arrow/.style={>=stealth, thick, <->}
}

\usefonttheme{professionalfonts}
\usefonttheme{serif}

\mode<presentation>
{
    \usenavigationsymbolstemplate{}
    \setbeamertemplate{headline}
    {
        \begin{beamercolorbox}[sep=1ex,wd=\paperwidth]{whitebox}
            \hfill{\color{awesome}\insertframenumber}\hspace*{0.4em}
        \end{beamercolorbox}
    }

    \setbeamercovered{invisible}
    \setbeamercolor{frametitle}{fg=awesome}
    \setbeamercolor{date in head/foot}{fg=awesome}
    \setbeamercolor{itemize item}{fg=awesome}
    \setbeamercolor{itemize subitem}{fg=awesome}
    \setbeamercolor{itemize subsubitem}{fg=awesome}
    \setbeamercolor{enumerate item}{fg=awesome}

    \setbeamertemplate{itemize item}[square]
    \setbeamertemplate{itemize subitem}[circle]
    \setbeamertemplate{itemize subsubitem}[triangle]
}

\setbeamertemplate{footline}[text line]{%
  \parbox{\linewidth}{\vspace*{-6pt}{\color{lightergrey}\copyright Mohammad Abdulaziz}\hfill\insertshortauthor\hfill{\color{awesome}}}}
\setbeamertemplate{navigation symbols}{}

\newcommand{\chictitle}[2]{
  \begin{frame}[plain,noframenumbering]
    \begin{center}
      \vspace{2.5em}
      {\color{awesome}\Large\textbf{#1}} \\
      \vspace{0.8em}
      \small{#2}
      \vspace{3em}
      \begin{center}
        {\color{awesome}\large Mohammad Abdulaziz}
      \end{center}
    \end{center}
  \end{frame}
}


\setbeamersize{text margin left=8pt,text margin right=8pt}

\title[Interactive Theorem Proving: Introduction]{Interactive Theorem Proving: Introduction}
\author[7CCSACTP]{Mohammad Abdulaziz}
\date{Week 2}

\begin{document}

\chictitle{Interactive Theorem Proving: Introduction}{Computer-Assisted Theorem Proving (7CCSACTP)}

\begin{frame}{This Video}
\newcommand{\intro}{1}
\only<\intro>{
\begin{itemize}
  \item What is interactive theorem proving?
  \item What are we doing over the next three lectures?
  \item What do I expect you to already know?
  \item Some Organisational notes
\end{itemize}
}

\note<\intro->{
\scriptsize
In this video I will
\begin{itemize}
  \item Introduce you to interactive theorem proving, the main idea, and motivate why I think it is an interesting technology to learn about
  \item I will then discuss briefly what are we doing over the next three lectures, which talk about ITPing, including the learning objectives
  \item I will also briefly discuss what concepts and techniques I expect you to already know from previous modules, although I will briefly introduce the different concepts throughout the lectures
  \item I will discuss some organisational notes regarding these next lectures and the coursework
 \end{itemize}
}
\end{frame}

\begin{frame}{What is ITP'ing?}
\newcommand{\syllogism}{1}
\newcommand{\formalproof}{2}
\newcommand{\axioms}{3}
\newcommand{\whyformal}{4}
\newcommand{\verybigproofs}{5}
\newcommand{\computerbasedformalproofs}{6}
\newcommand{\computerbasedformalproofseg}{7}
\newcommand{\formalproofmaths}{8}
\newcommand{\formalproofCS}{9}
\newcommand{\formalproofmywork}{10}

\note<\syllogism>{
\tiny
  \begin{itemize}
  \item Before I talk about interactive theorem proving I want to recapitulate with you the concept of formal proof
  \item Recall with me the following Aristotlian syllogism
    \begin{itemize}
    \item All men are mortal (Assume)
    \item Aristotle is a man (Assume)
    \item Aristotle is mortal (from 1, 2, and modus-ponens)
    \end{itemize}
  \item I hope we all agree that this is an unambiguously correct proof
  \end{itemize}
}

\only<\syllogism>{
  Formal proofs:
  \begin{itemize}
  \item Recall with me the following Aristotlian syllogism
    \begin{itemize}
    \item 1: All men are mortal (Assume)
    \item 2: Aristotle is a man (Assume)
    \item 3: Aristotle is mortal (from 1, 2, and modus-ponens)
    \end{itemize}
  \item Modus-ponens: if a implies b and a holds, then b holds
  \item I hope we all agree that this is an unambiguously correct proof
  \end{itemize}
}

\note<\formalproof>{
\tiny
  \begin{itemize}
  \item Our goal is to construct formal proofs of correctness for mathematical theorems and for algorithms/systems
  \item This is because formal proofs are the closest thing to objectively and unambiguously correct proofs:
    \begin{itemize}
    \item By that I mean, definitions and theorem statements whose meaning and well-formedness are unambiguously correct
    \item And also proofs that are unambiguously correct
    \item This is to be done wrt precisely and concisely stated axioms that indicate 
          what are correct reasoning or proof steps, and what is the meaning of different statements or claims
    \end{itemize}
  \end{itemize}
}

\only<\formalproof>{
  \begin{itemize}
  \item Our objective: construct formal proofs
  \item Unambiguous, correct proofs:
    \begin{itemize}
    \item Definitions' and theorem statements' meaning and well-formedness
    \item Proof correctness
    \end{itemize}
  \item All of that wrt precisely and concisely stated axioms
  \end{itemize}
}

\note<\axioms>{
  \tiny
  \begin{itemize}
  \item Proofs and definitions are written in formal language whose meaning is defined by set of axioms, e.g.
    \begin{itemize}
    \item Linear temporal logic axioms
    \item or Zermelo-Fraenkel set theory that are capable of representing any mathematical statement or thm, and they are more general than LTL. Some of its axioms are as follows
      \begin{itemize}
        \item Extensionality: {\small $\forall s,t.\; (\forall x.\; x\in s \leftrightarrow x\in t) \rightarrow s = t$}
        \item Regularity: {\small $\forall s.\; s\neq\emptyset \rightarrow \exists t.\; t\in s \wedge t\cap s =\emptyset$}
        \item Schema: {\small $\forall s,x_1,\dots,x_n.\exists t.\; \forall y.\; y\in t \leftrightarrow (y\in s \wedge \phi(s,x_1,\dots,x_n,y))$}
        \item \dots
      \end{itemize}
    \item there is also Type theory
    \item Peano Arithmetic
    \item First-order logic
    \item and Higher-order logic (HOL), which is the logical system of the theorem prover we use here
    \end{itemize}
  \end{itemize}
}

\only<\axioms>{
  \begin{itemize}
  \item Proofs and definitions are written in a formal language
  \item Meaning is defined by set of axioms, e.g.
    \begin{itemize}
    \item Linear temporal logic axioms
    \item More general: Zermelo-Fraenkel set theory:
      \begin{itemize}
        \item Extensionality: {\small $\forall s,t.\; (\forall x.\; x\in s \leftrightarrow x\in t) \rightarrow s = t$}
        \item Regularity: {\small $\forall s.\; s\neq\emptyset \rightarrow \exists t.\; t\in s \wedge t\cap s =\emptyset$}
        \item Schema: {\small $\forall s,x_1,\dots,x_n.\exists t.\; \forall y.\; y\in t \leftrightarrow (y\in s \wedge \phi(s,x_1,\dots,x_n,y))$}
        \item \dots
      \end{itemize}
    \item Other formalisms
      \begin{itemize}
      \item Type theory, Peano Arithmetic, First-order logic, and {\color{awesome}Higher-order logic (HOL)}
      \end{itemize}
    \end{itemize}
  \end{itemize}
}

\note<\whyformal>{
\tiny
  \begin{itemize}
  \item The initial reason to bother with formal proofs, which is still very relevant (later on that), is to ``formalise'' mathematics, which was more of an intellectual curiosity
    \begin{itemize}
    \item Some of the earliest trials to get formal proofs, i.e.\ to fully specify the axioms and show how every step in the proof follows from those axioms, are Euclid's elements (~300 B.C.) and principia mathematica [Russel and Whitehead 1910] 
    \end{itemize}
  \end{itemize}
}

\only<\whyformal>{
  \begin{itemize}
  \item Why formal proofs?
  \item Initially to ``formalise'' mathematics
    \begin{itemize}
    \item E.g.\ Euclid's elements (~300 B.C.) and principia mathematica [Russel and Whitehead 1910] 
    \end{itemize}
  \end{itemize}
}

\note<\verybigproofs>{
\tiny
  \begin{itemize}
  \item Then in the latter part of the 20th century it was needed for mathematics. Why?
  \item Proofs became too big to be reviewed by humans:
    \begin{itemize}
    \item Four-colour theorem [Appel and Haken 1976] (proof contains 1936 case splits)
    \item Classification of finite simple groups [100+ authors, 1955-2004] (proof is spread over ~500 journal articles, spanning 15k pages)
    \item Kepler's conjecture on Sphere packing [Hales 1998] (99\% correct, Ann.\ Math.)
    \end{itemize}
  \end{itemize}
}

\only<\verybigproofs>{
\begin{itemize}
  \item Then in the latter part of the 20th century it was needed for mathematics. Why?
  \item Proofs became too big to be reviewed by humans:
    \begin{itemize}
    \item Four-colour theorem [Appel and Haken 1976] (proof contains 1936 case splits)
    \item Classification of finite simple groups [100+ authors, 1955-2004] (proof is spread over ~500 journal articles, spanning 15k pages)
    \item Kepler's conjecture on Sphere packing [Hales 1998] (99\% correct, Ann.\ Math.)
    \end{itemize}
\end{itemize}
}

\note<\computerbasedformalproofs>{
\tiny
  \begin{itemize}
  \item This is what motivated computer-based formal proofs
  \item Idea: have a very simple computer program check every step in the proof, wrt to precisely stated axioms
  \item This program is called the \emph{logical kernel}
  \item Key property: \emph{the De-Brujin criterion}:
    \begin{itemize}
      \item Correctness of the whole formal proof depends ONLY on the correctness of the logical kernel!
      \item Because the logical kernel is very simple, we can have very high trust in the correctness of the proof
    \end{itemize}
  \end{itemize}
}

\only<\computerbasedformalproofs>{
\begin{itemize}
  \item Computer-based formal proofs:
    \begin{itemize}
    \item A computer program checks every step in the proof, wrt precisely stated axioms
    \item Logical kernel: very simple program performing the checking
    \item \emph{De-Brujin criterion}: correctness of the proof depends ONLY on the correctness of the logical kernel
    \end{itemize}
\end{itemize}
}

\note<\computerbasedformalproofseg>{
\tiny
  \begin{itemize}
  \item Some notable computer-based formal proofs in mathematics:
    \begin{itemize}
      \item Four-colour theorem in Coq [Gonthier 2005]
      \item Feit-Thompson theorem in Coq [Gonthier et al. 2012]
      \item Kepler's conjecture in Isabelle/HOL and HOL-Light [Hales et al. 2017]
      \item Liquid Tensor Experiment in Lean [Scholze, Commelin et al. 2021]
    \end{itemize}
  \end{itemize}
}

\only<\computerbasedformalproofseg>{
\begin{itemize}
  \item Computer-based formal proofs in mathematics:
    \begin{itemize}
      \item Four-colour theorem in Coq [Gonthier 2005]
      \item Feit-Thompson theorem in Coq [Gonthier et al. 2012]
      \item Kepler's conjecture in Isabelle/HOL and HOL-Light [Hales et al. 2017]
      \item Liquid Tensor Experiment in Lean [Scholze, Commelin et al. 2021]
    \end{itemize}
\end{itemize}
}

\note<\formalproofCS>{
\tiny
  \begin{itemize}
  \item In CS formal proof is used for \emph{formal verification}:
    \begin{itemize}
    \item A system and its desired properties are specified in a formal language
    \end{itemize}
  \item Notable applications in CS:
    \begin{itemize}
    \item seL4: formally verified OS microkernel [Klein et al. 2009]
    \item CompCert: verified compiler for C [Leroy 2009]
    \item CakeML: verified compiler for ML [Kumar et al. 2014]
    \end{itemize}
  \end{itemize}
}

\only<\formalproofCS>{
\begin{itemize}
  \item In CS formal proof is used for \emph{formal verification}:
    \begin{itemize}
    \item A system and its desired properties are specified in a formal language
    \end{itemize}
  \item Notable applications in CS:
    \begin{itemize}
    \item \textbf{seL4}: OS Microkernel [Klein et al. 2009]
    \item \textbf{CompCert}: compiler for C [Leroy 2009]
    \item \textbf{CakeML}: compiler for ML [Kumar et al. 2014]
    \end{itemize}
  \item Adopted by industry:
    \begin{itemize}
    \item Intel, Arm, AWS, Apple
    \end{itemize} 
  \end{itemize}  
}

\note<\formalproofmywork>{
\tiny
\begin{itemize}
  \item Research areas in interactive theorem proving:
  \begin{itemize}
    \item Verification of AI, algorithms, and software
    \item Formalisation of mathematics (complexity theory and combinatorial optimisation)
  \end{itemize}
\end{itemize}
}

\only<\formalproofmywork>{
\begin{itemize}
  \item Applications \& Research in ITP:
  \begin{itemize}
    \item Verification of AI, algorithms, and software
    \item Formalisation of mathematics (complexity theory \& combinatorial optimisation)
  \end{itemize}
\end{itemize}
}
\end{frame}

\begin{frame}{What will we do (i.e.\ learning outcomes)?}
\note{
\begin{itemize}
  \item In the next lectures I will introduce you to formal proofs in the context of the interactive theorem prover Isabelle/HOL
  \item I will focus on how to prove properties of simple functional programs
  \item I will also show you how to find such proofs on a notebook and then how you implement such proofs in Isabelle/HOL
  \item Take your own notes based on the videos!
\end{itemize}
}

\begin{itemize}
\item In these lectures:
  \begin{itemize}
    \item You will formally prove properties of simple functional programs
    \item You will do this pen-and-paper and using the ITP Isabelle/HOL
  \end{itemize}
\item Take your own notes from the videos!
\end{itemize}
\end{frame}

\begin{frame}{Expected Background Knowledge}
\note{
Although I will briefly re-introduce you to the following concepts, I expect you to have some familiarity with:
\begin{itemize}
  \item Mathematical induction
  \item Familiarity with first-order logic and how to formulate statements in it
  \item Familiarity with functional programming
\end{itemize}
}
\begin{itemize}
  \item Mathematical induction
  \item First-order logic and how to formulate statements in it
  \item Basic functional programming
  \item I will refresh them throughout the course!
\end{itemize}
\end{frame}

\begin{frame}{Organisational Points}
\newcommand{\downloadisabelle}{1}
\newcommand{\books}{2}
\newcommand{\assesment}{3}

\note<\downloadisabelle>{
\begin{itemize}
  \item As I mentioned earlier, we will be using the \href{https://isabelle.in.tum.de/}{\textcolor{awesome}{Isabelle/HOL}} theorem prover
  \item It was initially developed by Larry Paulson from the University of Cambridge, and is maintained by Cambridge, TU Munich, and many others
  \item Please download it from the link provided on this slide
  \item In addition to asking on the course forum, I recommend you sign up to \href{https://isabelle.zulipchat.com/}{\textcolor{awesome}{Isabelle's Zulip channel}}
\end{itemize}
}

\only<\downloadisabelle>{
\begin{itemize}
  \item We will use the \href{https://isabelle.in.tum.de/}{\textcolor{awesome}{Isabelle/HOL}} theorem prover
    \begin{itemize}
    \item Developed by Cambridge, TU Munich, and contributors worldwide
    \end{itemize}
  \item Questions: Course Forum + \href{https://isabelle.zulipchat.com/}{\textcolor{awesome}{Isabelle Zulip Community}}
\end{itemize}
}

\note<\books>{
\begin{itemize}
  \item Recommended Books:
  \item \href{http://concrete-semantics.org/}{\textcolor{awesome}{Concrete Semantics}} by Tobias Nipkow and Gerwin Klein
  \item \href{https://fdsa-book.net/}{\textcolor{awesome}{Functional Data Structures and Algorithms}} (Nipkow, Abdulaziz et al.)
  \item Brief Isabelle documentation: \href{https://isabelle.in.tum.de/doc/prog-prove.pdf}{\textcolor{awesome}{Programming and Proving}}
\end{itemize}
}

\only<\books>{
\begin{itemize}
  \item \href{http://concrete-semantics.org/}{\textcolor{awesome}{Concrete Semantics}} (Nipkow \& Klein)
  \item \href{https://fdsa-book.net/}{\textcolor{awesome}{Functional Data Structures and Algorithms}} (Nipkow, Abdulaziz et al.)
  \item Brief Isabelle documentation: \href{https://isabelle.in.tum.de/doc/prog-prove.pdf}{\textcolor{awesome}{Programming and Proving}}
\end{itemize}
}

\note<\assesment>{
\begin{itemize}
\item Assessment:
  \begin{itemize}
  \item 50\% Coursework (includes Viva)
  \item 50\% Exam (paper-based written examination)
  \end{itemize}
\end{itemize}
}

\only<\assesment>{
\begin{itemize}
\item Assessment:
  \begin{itemize}
  \item \textbf{50\% Coursework}: Formal proof projects \& Viva defense
  \item \textbf{50\% Exam}: Paper-based written examination
  \end{itemize}
\end{itemize}
}
\end{frame}

\end{document}
